
\begin{frame}[fragile]{Vues}

Toute requête produit une relation. Nommer cette requête
c'est nommer la relation résultat et pouvoir l'interroger.
\alert{Une vue est une requête nommée}.

\vfill

C'est aussi une table \alert{calculée} au moment où on l'interroge.

\vfill

Dans cette session:

\begin{redbullet}
  \item La commande \sqlkw{create view}
  \item Interrogation d'une vue
  \item Mise à jour d'une vue
\end{redbullet}

\vfill
\begin{blocgauche}
\alert{Ces diapositives correspondent au support en ligne disponible sur
le site \url{http://sql.bdpedia.fr/}}
\end{blocgauche}

\end{frame}

%%%%%%%%%%%%%%%%%%%%
\begin{frame}[fragile]{La commande \sqlkw{create view}}

Une vue est 
le résultat d'une requête 

\vfill

\begin{minted}[baselinestretch=.9,
               fontsize=\small,
               framesep=2mm]{sql}
create view Koudalou as 
        select nom, adresse, count(*) as nb_apparts 
        from Immeuble as i, Appart as a 
        where i.nom='Koudalou'
        and i.id=a.idImmeuble
        group by i.id, nom, adresse
\end{minted}

\vfill

\begin{blocgauche}
Le résultat de la requête
est réévalué à chaque fois que l'on accède à la vue. 
\end{blocgauche}
\end{frame}


%%%%%%%%%%%%%%%%%%%%
\begin{frame}[fragile]{Interrogation d'une vue}

On interroge une vue comme n'importe quelle table.

\vfill

\begin{minted}[baselinestretch=.9,
               fontsize=\small,
               framesep=2mm]{sql}
select * from Koudalou
\end{minted}


\vfill

\begin{tabular}{c|c|c}
\textbf{nom} & \textbf{adresse} & \textbf{nb\_apparts}\\ \hline 
Koudalou & 3 rue des Martyrs & 5\\ \hline 
\end{tabular}

\vfill

\begin{blocgauche}
Une vue peut répondre à des objectifs de \alert{simplification}
et/ou de \alert{confidentialité}.
\end{blocgauche}

\end{frame}



%%%%%%%%%%%%%%%%%%%%
\begin{frame}[fragile]{Vue dénormalisée}

Une vue peut présenter une base dénormalisée.

\vfill

\begin{minted}[baselinestretch=.9,
               fontsize=\small,
               framesep=2mm]{sql}
create view AppartKoudalou as 
       select no, surface, niveau, i.nom as immeuble, adresse,
              concat(p.prénom,  ' ', p.nom) as occupant 
       from Immeuble as i, Appart as a, Personne as p
       where i.id=a.idImmeuble
       and   a.id=p.idAppart
       and   i.nom='Koudalou'
\end{minted}

\vfill
\begin{small}
\begin{tabular}{c|c|c|c|c|c}
\textbf{no} & \textbf{surf.} & \textbf{niv.} & \textbf{immeuble} & \textbf{adresse} & \textbf{occupant}\\ \hline 
1 & 150 & 14 & Koudalou & 3 rue des Martyrs & Léonie Atchoum\\ \hline 
51 & 200 & 2 & Koudalou & 3 rue des Martyrs & Barnabé Simplet\\ \hline 
52 & 50 & 5 & Koudalou & 3 rue des Martyrs & Alice Grincheux\\ \hline 
43 & 75 & 3 & Koudalou & 3 rue des Martyrs & Brandon Timide\\ \hline 
\end{tabular}
\end{small}
\end{frame}


%%%%%%%%%%%%%%%%%%%%
\begin{frame}[fragile]{Insertion dans une vue}

Beaucoup de restrictions.


\vfill
\begin{redbullet}
  \item la vue doit être basée sur une seule table;
   \item toute colonne non référencée
    dans la vue doit pouvoir être mise à 
    \sqlkw{null} ou disposer d'une valeur par défaut;
  \item on ne peut pas mettre à jour un attribut
    qui résulte d'un calcul ou d'une opération.
\end{redbullet}

\vfill

\begin{minted}[baselinestretch=.9,
               fontsize=\small,
               framesep=2mm]{sql}
create view PropriétaireAlice
   as select * from Propriétaire
  where idPersonne=2

insert into PropriétaireAlice values (2, 100, 20)
insert into PropriétaireAlice values (3, 100, 20)
\end{minted}

\end{frame}


%%%%%%%%%%%%%%%%%%%%
\begin{frame}[fragile]{La clause \sqlkw{check option}}

Sur la vue précédente, la requête:

\begin{minted}[baselinestretch=.9,
               fontsize=\small,
               framesep=2mm]{sql}
select * from PropriétaireAlice
\end{minted}

ne montre pas le propriétaire 3 que l'on vient d'insérer !

\vfill

La clause \sqlkw{check option} permet de n'insérer que des nuplets
que l'on peut ensuite sélectionner.


\begin{minted}[baselinestretch=.9,
               fontsize=\small,
               framesep=2mm]{sql}
create view PropriétaireAlice
as select * from Propriétaire
where idPersonne=2
with check option
\end{minted}

\end{frame}

%%%%%%%%%%%%%
\begin{frame}{À retenir}
Les \alert{vues} sont des requêtes nommées que l'on peut traiter
comme des tables.

\vfill

Elles permettent de restructurer ``virtuellement''  une base.

\begin{redbullet}
   \item  Pour faciliter l'accès (jointures pré-définies)
   \item Pour restreindre la visibilité des données (on ne donne accès qu'à la vue)
   \item Les insertions dans les vues offrent peu d'intérêt.
\end{redbullet}


\end{frame}

